% ============================================================================
% Response to Reviewers - TSAS revision round 1
% "Tesseract: Generating Indoor Graph Maps from Large-Scale Low-Semantic
%  Floor Plans"
%
% Conventions:
%   - Reviewer comments are quoted VERBATIM in that reviewer's color.
%   - Responses are in black and every response opens by thanking the
%     reviewer for the specific point.
%   - Every change is localized by section, subsection, and figure or
%     table number. Page numbers [p. XX] are filled after the final
%     compile of the revised manuscript.
%   - The same colors mark the corresponding edits in main_revision_1.tex.
% ============================================================================
\documentclass[11pt]{article}
\usepackage[margin=1in]{geometry}
\usepackage{xcolor}
\usepackage{booktabs}
\usepackage{longtable}
\usepackage[normalem]{ulem}
\usepackage{parskip}

\definecolor{RevOneCol}{RGB}{0,84,159}     % Reviewer 1 - blue
\definecolor{RevTwoCol}{RGB}{0,122,94}     % Reviewer 2 - green
\definecolor{RevThreeCol}{RGB}{118,42,131} % Reviewer 3 - purple
\definecolor{RevAllCol}{RGB}{178,58,0}     % Editor / multiple - vermilion

\newcommand{\commentone}[2]{\par\noindent\textbf{Comment #1.}\begin{quote}\itshape\color{RevOneCol}#2\end{quote}}
\newcommand{\commenttwo}[2]{\par\noindent\textbf{Comment #1.}\begin{quote}\itshape\color{RevTwoCol}#2\end{quote}}
\newcommand{\commentthree}[2]{\par\noindent\textbf{Comment #1.}\begin{quote}\itshape\color{RevThreeCol}#2\end{quote}}
\newcommand{\resp}{\par\noindent\textbf{Response.}~}
\newcommand{\loc}[1]{\par\noindent\fbox{\parbox{0.97\linewidth}{\small\textbf{Where in the revision:} #1}}\par\medskip}

\title{Response to Reviewers \\ \large Manuscript: ``Tesseract: Generating Indoor Graph Maps from Large-Scale Low-Semantic Floor Plans'' (ACM TSAS)}
\author{Yaqoob Ansari \and Eduardo Feo Flushing \and Khaled A. Harras}
\date{Revision 1, [DATE]}

\begin{document}
\maketitle

\noindent Dear Editor-in-Chief, Associate Editor, and Reviewers,

We are grateful for the time and care invested in these reviews. The three reports identified the places where the submission fell short of its own claims, most notably the degree of multi-floor automation, the precision of the funneling guarantee, and the breadth and objectivity of the evaluation. Addressing them has changed the system and not only the paper. The revision adds an automatic inter-floor transition matching module that replaces the manual mapping as the default path, a formally stated and proven connectivity guarantee, and eleven new or extended experiments including degradation studies, ablations, an objective route-fidelity reference, and a quantitative baseline comparison. We believe the manuscript is substantially stronger for it, and we thank the reviewers for pushing it there.

Every change in the revised manuscript is color-coded by the reviewer who motivated it:

\begin{center}
\begin{tabular}{ll}
\textcolor{RevOneCol}{\rule{2em}{0.8em}} & Reviewer 1 (blue) \\
\textcolor{RevTwoCol}{\rule{2em}{0.8em}} & Reviewer 2 (green) \\
\textcolor{RevThreeCol}{\rule{2em}{0.8em}} & Reviewer 3 (purple) \\
\textcolor{RevAllCol}{\rule{2em}{0.8em}} & Raised by multiple reviewers (vermilion) \\
\end{tabular}
\end{center}

A summary table of changes follows, and detailed point-by-point responses are below. Each response is localized by section, subsection, and figure or table number in the revised manuscript.

% Table of changes - \input into response_to_reviewers.tex
\section*{Summary of Changes}
\begin{longtable}{p{0.5cm} p{2.0cm} p{4.6cm} p{4.9cm} p{2.9cm}}
\toprule
\# & Reviewer(s) & Concern & Change in revision & Location \\
\midrule
\endhead
1 & \textcolor{RevAllCol}{R1.9, W1, R3.2} & Multi-floor step manual, delta weak, automation overstated & Automatic inter-floor transition matching, registration plus optimal assignment with rejection, manual mapping demoted to override, claims re-audited & Sec.\ 9.2, 10.8, Figs.\ 7--8 \\
2 & \textcolor{RevOneCol}{R1.4} & Algorithm 1 does not guarantee main-node routing & Guarantee corrected and proven, Algorithm 1 rewritten, bypass behavior stated with measurement, three-way ablation added & Sec.\ 7.3, 10.4, Tab.\ 3 \\
3 & \textcolor{RevAllCol}{R2.W3, R3.3} & Corridor-door failure under-treated, projection unsupported & Monte-Carlo knockout curves with per-door sensitivity, bridge doors identified as cut vertices, repair direction stated & Sec.\ 10.3, 11, Fig.\ 3 \\
4 & \textcolor{RevAllCol}{R1.7, W7, R3.5} & Euclidean fidelity proxy weak, nominal transition weights & Free-space geodesic reference with dual bounds and verification overlays, inter-floor edges re-weighted by modeled vertical travel & Sec.\ 10.6, 9.3, Fig.\ 6 \\
5 & \textcolor{RevOneCol}{R1.1} & Applications and corridor density unjustified & Downstream-task validation plus measured density trade-off & Sec.\ 10.9, 10.5, Fig.\ 5 \\
6 & \textcolor{RevAllCol}{R1.6, W2, R3.1c} & Evaluation narrowed to one site & Exhaustive metric layer over all plans, dataset table, conference sites being restored & Sec.\ 10.1, 10.4, Tab.\ 2 \\
7 & \textcolor{RevThreeCol}{R3.1c} & No baseline comparison & Quantitative medial-axis baseline, decisively outperformed on every axis & Sec.\ 10.7, Tab.\ 4 \\
8 & \textcolor{RevTwoCol}{R2.W6} & Missing ablations & Funneling, keep-term, and pruning-cost ablations plus matcher provenance & Sec.\ 10.4, 9.2, Tab.\ 3 \\
9 & \textcolor{RevOneCol}{R1.8} & Text-detection accuracy unreported & Measured detection and recognition with failure-consequence analysis & Sec.\ 10.2 \\
10 & \textcolor{RevOneCol}{R1.3, R1.17} & Label-as-main-node unjustified, term undefined & Placement measurement over 172 rooms, term removed & Sec.\ 6.3, 10.2 \\
11 & \textcolor{RevTwoCol}{R2.W4, W5} & Door semantics and corridor identification unclear & Precise definitions, vocabulary table with fuzzy matching, label-driven semantics corrected, failure demonstrated & Sec.\ 3.1, 3.2, 5.2, Tab.\ 1 \\
12 & \textcolor{RevThreeCol}{R3.4} & Applicability overclaimed & Operating-envelope statement with exclusions and demonstrated failure mode & Sec.\ 3.1, 11 \\
13 & \textcolor{RevOneCol}{R1.5} & Show the floor plans & Dataset figure plus full-width appendix reproductions & Sec.\ 10.1, App.\ B, Fig.\ 2 \\
14 & \textcolor{RevOneCol}{R1.2, minors} & Representation, acronyms, OCR mechanism, figure fixes, paragraph placement & Each addressed individually in the point-by-point responses R1.2 and R1.10 to R1.17 & Sec.\ 3.1, 5.2, 6.1, 8.3, 10.4 \\
\bottomrule
\end{longtable}


\clearpage
% ============================================================================
\section*{Reviewer 1}
% ============================================================================

We sincerely thank Reviewer 1 for an unusually careful reading. The review caught a genuine technical flaw in Algorithm 1 that we had missed ourselves, and its broader questions about applications, evaluation breadth, and the multi-floor delta shaped the largest additions of this revision.

\commentone{R1.1}{Most importantly maybe, what exactly are your targeted applications? For example, why would you need / want corridors to be densely covered with nodes (and edges), at least in the original graph? For human navigation, for example, it appears sufficient to have a node at every `entry' to the corridor. And what is different in / with corridors than rooms that rooms do not receive the same treatment?}
\resp We thank the reviewer for pressing on this, because the question exposed that our design choices were asserted rather than justified, and answering it produced two new studies. First, the revision validates concrete applications directly instead of listing them. A new downstream-task subsection demonstrates step-free accessibility routing, which is expressible only because transitions are typed, and per-room egress-distance analysis. Second, we ran a corridor-density ablation with full pipeline reruns at grid spacings from 20 to 80 pixels. The result answers the reviewer's intuition precisely. A 60-pixel grid keeps the graph fully connected with room-to-exit routes within 4.6\% of the 20-pixel default at 57\% fewer nodes, so dense coverage is indeed not required for reachability. At 80 pixels, however, the corridor mesh fragments catastrophically and completeness collapses to 0.49. The failure mode is a cliff rather than a slope, and the dense default buys a safety margin against narrow corridors at a now-quantified cost. On the final question, rooms and corridors both receive spatial subnodes, and the difference is the seeding. Rooms are compact labeled regions grown by flood fill and sampled hexagonally, whereas corridors are elongated unlabeled regions sampled on a uniform grid over the corridor mask. The revised text states this rationale explicitly.
\loc{Applications validated in Section 10.9 (Downstream Task Validation) [p.~XX]. Density ablation in Section 10.5 (Corridor Sampling Density) with Figure 5 [p.~XX]. Corridor-versus-room sampling rationale in Section 6.4 (Corridor and Outdoor Sampling) [p.~XX]. All marked in blue.}

\commentone{R1.2}{Why do you perform spatial segmentation on a pixel level? Because you (still) operate on the image here? In principle, the extracted walls could provide boundaries of `closed' spaces with everything inside being part of that space, without explicitly assigning a `room' value to every pixel. That is, in other words, what is the underlying representation here and for subsequent steps? Is it the image space with pixels as the smallest unit?}
\resp We thank the reviewer for asking us to make the representation explicit, which the submission never did. The answer is yes. The underlying representation throughout the pipeline is raster image space, and the pixel is the atomic unit of every stage. The revision states this at the start of the problem setting. We also now give the design rationale the reviewer invites. We fill at pixel granularity rather than polygonizing wall boundaries because legacy drawings frequently contain hairline gaps, decorative strokes, and symbol overlaps that make closed-polygon extraction brittle, whereas flood fill degrades gracefully under all three. Vectorized boundaries are noted as attractive future work.
\loc{Representation statement opens Section 3.1 (Input Assumptions) [p.~XX]. Design rationale in Section 6.2 (Flood Fill for Room Segmentation) [p.~XX]. Marked in blue.}

\commentone{R1.3}{Why is it reasonable to use the position of the text label for a room as the position for the main node? The label could be arbitrarily placed in principle. What makes that location so special?}
\resp We thank the reviewer for this question, because the honest answer is that nothing forces labels to be central, so instead of arguing we measured it. Across 172 rooms on nine processed floor sections, the offset between the label position and the geometric centroid of the room's flood-filled region has a median of 0.24 room-equivalent radii, a mean of 0.27, and a worst case of 0.98, and 92\% of labels fall inside the room's inscribed circle. Architectural drafting practice does place labels centrally, and the revision reports this measurement as the empirical justification rather than assuming it.
\loc{Measurement summarized in Section 6.3 (Subnode Placement) and reported in Section 10.2 (Text and Semantics) [p.~XX]. Marked in blue. See also R1.17 for the removed terminology.}

\commentone{R1.4}{And then (page 9, `room family funneling') you state that you root all `subnode-to-door' paths through the main node. I do fail to see how Algorithm 1 actually guarantees this. You seem to compute all shortest paths from any node in the room to the one closest to the door. But these paths do not necessarily contain the main node. All you do (as far as I can tell) is add an edge from the main node to the door as well. You may well end up with `subnode-to-door' paths that do not go through the main node. Or am I missing something? Along the same lines: why would you want all paths in a room to go through the main node?}
\resp We are genuinely grateful for this observation. The reviewer is not missing anything. The analysis is exactly right, and we verified it in our own implementation, where the closest-family-node edge and the main-node edge are both retained, so that on our primary evaluation floor 69 of 186 subnode-to-door shortest paths bypass the main node. The submitted phrasing overclaimed, and we have replaced it with the property that actually holds and that downstream routing actually needs. The revision rewrites Algorithm 1 to describe the real construction, a transient proximity lattice with an anchor-rooted shortest-path tree and a shortcut edge per door. A new Proposition with proof establishes complete intra-room connectivity, meaning every family node reaches every room door, at linear intra-room edge cost, since the retained edges form a spanning tree of at most $|S|$ edges versus $\Theta(|S| \cdot |D_r|)$ for dense wiring. The bypass behavior is now stated openly as intentional, with the measured 69 of 186 figure. On the reviewer's final question, we agree entirely, and the revised text says so. We do not want all paths through the main node, because the shortcut edge strictly improves realized path lengths. A three-way ablation quantifies the trade-off. Funneling stays complete at 167 intra-room edges versus 1,869 for dense wiring with only 7.2\% route stretch, and a second floor yields exactly one intra-room edge per subnode, matching the Proposition. We believe the reviewer's catch converted a weak claim into one of the strongest results in the paper, and we thank them for it.
\loc{Rewritten Algorithm 1, Proposition 1 with proof, the correction paragraph, and the complexity analysis are all in Section 7.3 (Room-Family Funneling) [p.~XX]. The ablation is in Section 10.4 (Graph Quality, Ablations, and Pruning) with Table 3 [p.~XX]. Marked in blue, with the shared ablation passages in green.}

\commentone{R1.5}{I strongly encourage (almost demand ;) the authors to include a figure with the floor plans that they use in the evaluation. As it stands, it is really hard to tell how `easy' or `complex' the task is [...] I'd be perfectly fine with having these figures in an appendix if they take up too much space in the main text.}
\resp We thank the reviewer for this practical suggestion and fully agree that readers need to see the inputs. The revision adds a dataset figure with two representative drawings and a caption that spells out the annotation conventions, and the appendix reproduces the drawings at full page width so labels, door symbols, and wall conventions are legible.
\loc{Figure 2 in Section 10.1 (Dataset and Setup) [p.~XX] and full-width reproductions in Appendix B (Additional Visualizations) [p.~XX]. Marked in blue.}

\commentone{R1.6}{More importantly, having just a single example / set of floor plans makes for a somewhat poor evaluation. [...] I'd strongly encourage you to include two or three additional buildings in the evaluation, ideally not from the same agency [...] By the way, the SIGSpatial paper seems to include three sites. Why are they not all included in this paper any longer?}
\resp We thank the reviewer for holding the journal version to at least the conference standard, and we agree the narrowing was a mistake. The revision broadens the evaluation on two axes. First, the graph-quality metrics, now computed exhaustively from full component structure by a dedicated metric layer rather than sampled, are reported for all processed plans, including ten stylistically distinct smaller plans, and the dataset table lists the corpus explicitly. Second, the conference sites are being restored so that the journal version strictly extends the conference evaluation. The omission was one of scope rather than capability, and we appreciate the reviewer flagging it. [The consolidated multi-site results are marked in the manuscript and will be finalized in this revision.]
\loc{Dataset table is Table 2 in Section 10.1 (Dataset and Setup) [p.~XX]. All-plans metrics in Section 10.4 (Graph Quality, Ablations, and Pruning) [p.~XX]. Marked in vermilion, jointly with R2.W2 and R3.1c.}

\commentone{R1.7}{I understand the aim to maintain `true distances.' However, even on pixel-level precision, the navigable shortest path ($d_{graph}$) between two rooms will in almost all circumstances be (substantially) longer than the direct line distance between the two centroids [...] Why is this a meaningful metric then? However, later on in Section 10.7 you claim that $d_{graph}$ is very close to $d_{euclidean}$. [...] Please clarify!}
\resp We thank the reviewer for this analysis, which we found convincing enough that we did not merely clarify the metric but replaced it. The revision validates route quality against the free-space geodesic, the exact shortest walkable path through the drawing's non-wall pixels with doorways opened at detected door positions, computed on the raster. This reference respects walls, which straight-line distance does not. Against it, graph routes on the primary enclosed floor are within 4.7\% of optimal on average with a ninetieth percentile of 9.3\%. We report the reference as two bounds that bracket the truth, a permissive bound over all free space and a conservative interior-only bound, and the gap between them on open sections is itself informative, since it quantifies sparse outdoor routing. Finally, we measured exactly the effect the reviewer described. The ratio of geodesic to Euclidean distance runs between 1.35 and 1.67 on average and reaches 3.3, which is the quantified reason the Euclidean proxy was misleading. Rendered overlays showing the geodesic and the graph route on the floorplan accompany the revision for direct visual verification.
\loc{Section 10.6 (Route Fidelity against Free-Space Geodesics) with Figure 6 [p.~XX]. Marked in vermilion, jointly with R2.W7 and R3.5.}

\commentone{R1.8}{You state that the major issues in the workflow arise from failed door detection. I can perfectly accept this. However, I'd still be interested to learn how good the text label detection works? Do you manage to extract all labels? 95\%? 80\%?... What happens if there are labels that are not identified (at all)?}
\resp We thank the reviewer for requesting numbers where we had offered none. The revision measures label extraction against a hand-enumerated ground truth of every label visible on the drawing. Detection finds 43 of 43 labels and recognition is exact for 41 of 43, or 95.3\%, with both failures truncating the leading digit of a room number, for example 1011 read as 101. We also answer the consequence question by failure class. A misrecognized numeric label still yields a valid room node whose name is wrong, so connectivity is unaffected. A label missed entirely leaves its region without a semantic seed and the region is not instantiated, which we demonstrate concretely with an unlabeled variant of the same drawing that produces zero transition nodes and fails the semantic-class requirement of door classification. In addition, the vocabulary matcher now carries single-edit fuzzy tolerance, so a single-character confusion in a semantic word, such as nna for na, resolves to the correct class instead of leaking into the room set.
\loc{Section 10.2 (Text and Semantics) [p.~XX], with the matching mechanism in Section 5.2 (Text Recognition) and the vocabulary in Table 1, Section 3.1 [p.~XX]. Marked in blue.}

\commentone{R1.9}{Finally, the only actual extension compared to the original paper seems to be the ability to combine several floor plans to a multi-level graph [...] the way this step is realized, it appears to be almost trivial from a scientific or algorithmic perspective [...] all that is left to do then is to add indices to all labels, and to perform some simple assertions. [...] to me makes for a very weak extension of the existing paper when it comes to judging the `delta' between the two manuscripts.}
\resp We thank the reviewer for stating this plainly, because the criticism was fair and it motivated the largest addition of the revision. The correspondence problem, deciding which stairs and elevators on adjacent floors are the same physical shaft, is now solved automatically rather than supplied by hand. The revised Section 9 formulates this in two stages. Registration estimates a similarity transform between adjacent floor drawings from four hypothesis generators, including phase correlation of glyph-filtered wall masks and a drawing-scale prior bootstrapped from detected door widths, ranked by a symmetric coverage-gated chamfer distance and refined by least squares and iterative closest point. Matching then solves a global optimal assignment with an explicit rejection option, so shafts that terminate remain unmatched rather than being force-paired, which no greedy nearest-shaft heuristic can offer. The hand-written mapping file is demoted to an optional override, and the validation rules now verify the machine's answer rather than a human's input. The robustness evaluation is systematic. Under known synthetic perturbations of a real floor, matching accuracy is 1.000 through 400 pixels of translation, 30 degrees of rotation, scale changes from 0.7 to 1.4, twin-shaft separations down to 25 pixels, and shaft-deletion tests, while a greedy nearest-neighbor baseline degrades to 0.75 on close twins, and registration recovery error stays between 0.5 and 4.9 pixels median. A cross-floor case study on two genuinely different drawings shows the intended conservative behavior, one supported correspondence accepted and every other shaft rejected with explicit warnings. [Real multi-floor precision and recall against the hand-mapping ground truth will be added with the consolidated dataset.] We hope the reviewer agrees that the delta now contains a substantive algorithmic contribution, joint registration and assignment with rejection, rather than bookkeeping plus assertions.
\loc{Section 9.2 (Automatic Transition Matching) with the Registration, Assignment, and Provenance paragraphs [p.~XX]. Robustness evaluation and the cross-floor case study in Section 10.8 (Multi-Floor Evaluation) with Figures 7 and 8 [p.~XX]. Marked in vermilion, jointly with R2.W1 and R3.2.}

\subsection*{Reviewer 1, minor comments}

\commentone{R1.10}{Please make sure to properly introduce all acronyms [...] for VGG-BiLSTM-CTC only the CTC part is spelled out / introduced.}
\resp We thank the reviewer for the careful catch. All acronyms are now introduced at first use, including VGG, BiLSTM, CTC, CRAFT, BIM, and SLAM.
\loc{Section 5.2 (Text Recognition) for the recognizer acronyms and Section 1 for BIM and SLAM [p.~XX]. Marked in blue.}

\commentone{R1.11}{Same paragraph: you state that you handle common OCR errors in text recognition. How do you `know' whether the system mixed up 0 and O etc.? That is, what mechanism is used here?}
\resp We thank the reviewer for asking for the mechanism rather than accepting the claim. The revision describes exactly what is implemented and nothing more. Recognized tokens are normalized by semantic context. Tokens are matched against the closed label vocabulary with case folding and single-edit fuzzy tolerance, so an alphabetic token of three or more characters matches a vocabulary word within Levenshtein distance one, ambiguous matches across two classes are rejected, and numeric tokens are excluded from fuzzy matching. Room identifiers are expected to be numeric, so alphabetic look-alikes inside otherwise numeric tokens are mapped to digits, which is how a letter O is resolved to the digit zero. The vaguer general claim has been removed.
\loc{Section 5.2 (Text Recognition), second paragraph [p.~XX]. Marked in blue.}

\commentone{R1.12}{In the Wall Mask Extraction part, do you `take out' the previously identified text from the image? I would assume that this text is usually also in dark / black color, so could be mistaken for a `wall' in principle.}
\resp We thank the reviewer for this sharp observation, and the revision answers it honestly in both directions with a measurement. In the segmentation pipeline, text strokes do enter the wall mask, and we measured them at roughly 9\% of wall ink on the primary evaluation floor. This contamination is harmless there because flood fill treats glyphs as small obstacles and the dilation buffer dominates. In the new registration module, where wall masks serve as alignment evidence and glyph noise would genuinely hurt, glyphs are removed explicitly by connected-component filtering before scoring. Both facts are now stated where each is relevant.
\loc{Section 6.1 (Wall Mask Extraction) [p.~XX] and Section 9.2 (Automatic Transition Matching, Registration paragraph) [p.~XX]. Marked in blue.}

\commentone{R1.13}{Page 8, Figure 8: [...] how do you get from subfigure 3 to subfigure 4? For most rooms, the main node (the red star) doesn't seem to be present anymore [...] How does that happen if you just remove nodes, but do not add new ones?}
\resp We thank the reviewer for catching this inconsistency, which we traced to the subfigures having been rendered from different runs. The figure is regenerated from a single run with consistent node identity across panels, and the caption now states the exact operation between each pair of panels.
\loc{The node-placement figure in Section 6 [p.~XX], regenerated. Marked in blue.}

\commentone{R1.14}{`Nodes with degree zero after edge removal are subsequently pruned, except for transition nodes[...]' Is it even possible for transition nodes to have a node degree of zero? Since you include shortest paths from rooms to these transition nodes, that shouldn't happen, right?}
\resp We thank the reviewer for probing the necessity of this protection. It is possible, exactly when no room-to-transition path exists before pruning, for instance when a missed door seals a shaft off from the corridor network. The revised text explains the case, and the new keep-term ablation demonstrates that the protection is load-bearing, since re-deriving the keep-set without the transition term silently removes every transition from the pruned graph.
\loc{Section 8.3 (Graph Pruning) [p.~XX] and the pruning ablation in Section 10.4 [p.~XX]. Marked in blue and green respectively.}

\commentone{R1.15}{Lines 517-520: ``This conservative approach...'' This paragraph should be part of the results / evaluation. Not part of the methods description.}
\resp We thank the reviewer for the structural point and agree. The passage is moved into the evaluation, where it now closes the pruning analysis with the measured justification.
\loc{Section 10.4 (Graph Quality, Ablations, and Pruning), end of the pruning paragraph [p.~XX]. Marked in blue.}

\commentone{R1.16}{`Corridor-to-corridor doors exhibit a bimodal distribution.' That may well be. The problem is that you cannot see / represent bimodal distributions in a box plot (Figure 3) as box plots obfuscate the properties of bimodal distributions.}
\resp We thank the reviewer for the statistical care. The box plot is replaced with a representation that shows distribution shape, and the bimodality claim is retained only where the new plot visibly supports it.
\loc{The door-statistics figure in Section 10 [p.~XX], replaced. Marked in blue.}

\commentone{R1.17}{Section 10.4: Why are the main node rooms positioned at `semantic centroids?' Or rather, what are `semantic centroids?' What makes them semantic?}
\resp We thank the reviewer for pressing on the term, which was imprecise. It is removed throughout. The accurate description, now used everywhere, is the label position, which the measurement of R1.3 shows to be empirically near-central with a median offset of 0.24 equivalent radii.
\loc{Terminology corrected throughout, with the measurement in Sections 6.3 and 10.2 [p.~XX]. Marked in blue.}

\clearpage
% ============================================================================
\section*{Reviewer 2}
% ============================================================================

We sincerely thank Reviewer 2 for the encouraging assessment of the problem relevance, the pipeline decomposition, and the room-family mechanism, and equally for the seven precisely stated weaknesses. Each one is addressed below, several with new experiments.

\commenttwo{R2.W1}{One of the main new contributions is the multi-floor connectivity framework, but this component appears to rely on structured transition mappings between stairs/elevators across floors. If these mappings are manually provided, and if floor identity is inferred from filename prefixes, the system is metadata-assisted rather than fully automatic. The authors should clarify the degree of automation.}
\resp We thank the reviewer for the exact characterization, which was accurate for the submission. Rather than only clarifying the wording, the revision changes the system. Inter-floor correspondences are now inferred automatically through registration and optimal assignment with rejection, described fully under R1.9, and the manual mapping file is retained only as an optional override. The revised text also states precisely what remains metadata-assisted by design, namely floor indices read from filename prefixes, and the automation claims in the abstract, introduction, and Section 9 have been re-audited to match the implementation exactly.
\loc{Section 9.2 (Automatic Transition Matching) [p.~XX], with tempered claims in the Abstract and Section 1 [p.~XX]. Marked in vermilion, jointly with R1.9 and R3.2.}

\commenttwo{R2.W2}{The journal evaluation is deeper for the multi-floor case but narrower in dataset diversity. The previous work evaluated multiple datasets, including Site A, Site B, and SESYD, whereas the journal version focuses mainly on one three-floor university building. This weakens claims about generalization.}
\resp We thank the reviewer for the fair comparison with the conference scope, and we agree. As detailed under R1.6, the metric layer now reports all processed plans including the ten stylistically distinct smaller plans, and the conference sites are being restored so the journal evaluation strictly extends the conference one. [The consolidated multi-site results are marked in the manuscript and will be finalized in this revision.]
\loc{Section 10.1 (Dataset and Setup, Table 2) and Section 10.4 [p.~XX]. Marked in vermilion, jointly with R1.6 and R3.1c.}

\commenttwo{R2.W3}{Door detection remains the dominant source of graph failure. In particular, corridor-to-corridor doors are much less reliable than room-to-corridor and exit doors. Since missed corridor doors can fragment the graph, this limitation deserves stronger treatment and possibly additional technical improvement.}
\resp We thank the reviewer for insisting on stronger treatment, because the resulting study both deepened and sharpened our own understanding of the failure mode. The revision adds a Monte-Carlo door-knockout analysis, deleting up to half of the detected doors over twenty draws per level, both at random and restricted by type, together with a per-door sensitivity analysis. The findings are more nuanced than the submission's claim. Most corridor-to-corridor doors turn out to be redundant because the corridor mesh reconnects around them. The danger concentrates in bridge doors, whose removal disconnects entire sections. On one evaluation floor a single such door costs 0.51 of pairwise completeness and 0.54 of exit reachability by itself, while room-to-corridor deletions cause bounded linear damage. Crucially, bridge doors coincide with the cut vertices of the graph, so the vulnerable doors are mechanically identifiable. This turns the dominant failure mode into a targeted verification and repair problem, and the discussion now outlines the concrete repair direction, verifying the wall image along candidate bridges and inserting flagged doors where a door-width opening exists, which we scope as immediate future work rather than claiming it.
\loc{Section 10.3 (Connectivity under Door-Detector Degradation) with Figure 3 [p.~XX], and the repair direction in Section 11 (Discussion, Door Detection and Repair) [p.~XX]. Marked in vermilion, jointly with R3.3.}

\commenttwo{R2.W4}{The paper should clarify whether doors are modeled as undirected or directed transitions, especially for exit doors, emergency exits, or one-way access constraints. The meaning of ``exit door'' should also be defined more precisely.}
\resp We thank the reviewer for requesting the precise semantics. The revision states them in the schema. The graph is undirected, and an exit door is defined as a door whose two incident regions are an interior space and an outdoor region. One-way restrictions and emergency-only egress are not modeled, this limitation is now listed explicitly, and directional edge attributes are identified as the natural extension.
\loc{Section 3.2 (Output Graph Schema) [p.~XX] and Section 11 (Discussion, Directional Constraints) [p.~XX]. Marked in green.}

\commenttwo{R2.W5}{The corridor identification process should be clarified. [...] it is unclear how the system handles unlabeled corridors, whether corridor regions can be inferred without explicit text labels, and whether color-coded corridor information in the evaluation data is used or ignored.}
\resp We thank the reviewer for the three-part question, which uncovered an inconsistency between the submitted text and the implementation that we are glad to correct. Semantics are label-driven, and the input-assumption text that suggested color-based identification has been corrected accordingly. Corridor regions are seeded by text labels such as hall and its synonyms, matched with case folding and single-edit fuzzy tolerance, and the full vocabulary with per-class requirements is now a table. Unlabeled corridors are not inferred. A drawing without corridor labels fails the semantic-class requirement and is rejected rather than silently mis-segmented, and we demonstrate this behavior on an unlabeled variant of an evaluation drawing.
\loc{Corrected input assumptions and Table 1 in Section 3.1 [p.~XX], matching mechanism in Section 5.2 [p.~XX], and the demonstrated failure in Section 10.2 [p.~XX]. Marked in green.}

\commenttwo{R2.W6}{The evaluation lacks sufficient ablation studies. It would be useful to isolate the contribution of the new room-family funneling mechanism, transition-node preservation, graph pruning, and multi-floor validation constraints.}
\resp We thank the reviewer for the specific list, and the revision adds one study per requested component. First, funneling is compared against dense subnode-to-door wiring and a single-node-per-room encoding on completeness, graph size, route stretch, and query time, and it occupies the size-quality sweet spot, complete at 167 intra-room edges versus 1,869 dense with 7.2\% stretch. Second, a keep-term ablation isolates transition preservation, showing that disabling the transition term silently removes every transition and disabling the exit term removes every exit. Third, the pruning study measures the full pass at a 59\% node and 82\% edge reduction with exactly zero route inflation and a 4.2 times query-time speedup. Fourth, the multi-floor validation constraints are exercised by the new automatic matcher, whose per-run provenance reports registration-quality warnings for weak wall agreement, single supporting pairs, and mirroring.
\loc{Funneling ablation with Table 3 and the pruning ablation in Section 10.4 [p.~XX]. Matcher provenance in Section 9.2 (Provenance and verification) [p.~XX]. Marked in green.}

\commenttwo{R2.W7}{The geometric fidelity metric is useful but limited. Comparing graph path length with Euclidean distance does not fully capture the quality of indoor routes [...] A comparison against manually annotated navigable paths would be more convincing.}
\resp We thank the reviewer for the push toward a navigationally realistic reference, and we adopted a stronger one than manual annotation. As detailed under R1.7, fidelity is now measured against exact free-space geodesics, which respect walls and doorways by construction, with graph routes within 4.7\% of optimal on the enclosed evaluation floor. Rendered per-pair overlays make the reference paths directly verifiable by inspection, which we believe delivers the convincingness the reviewer asks for with an objective and reproducible reference.
\loc{Section 10.6 (Route Fidelity against Free-Space Geodesics) with Figure 6 [p.~XX]. Marked in vermilion, jointly with R1.7 and R3.5.}

\clearpage
% ============================================================================
\section*{Reviewer 3}
% ============================================================================

We sincerely thank Reviewer 3 for recognizing the practical contribution and for laying out a precise improvement program. We followed it closely, and several of the revision's largest changes trace directly to this review.

\commentthree{R3.1a}{[The paper] should clearly position the extension around multi-floor graph construction and room-family funneling, rather than presenting door classification as a new contribution and claiming a comprehensive experimental evaluation. [...] Multi-floor graph construction and validated transition mapping together with room-family representation with distance-transform subnodes and funneling (i.e., contributions \#1 and \#2) are new substantial contributions.}
\resp We thank the reviewer for the clear editorial guidance and adopt it as stated. The contribution list is rewritten around the two substantial contributions, multi-floor graph construction, which now includes automatic transition correspondence, and the room-family representation with its proven funneling guarantee. Door classification is repositioned as a supporting deepening rather than a novelty.
\loc{Section 1 (Introduction), contribution paragraph and list [p.~XX]. Marked in purple.}

\commentthree{R3.1b}{Type-aware door classification (i.e., contribution \#3) already existed in the conference paper (i.e., Algorithm 2) and only the evaluation/treatment is expanded. Authors should clarify if the same algorithm is used in both works. If yes, then contribution \#3 should be rephrased to highlight the broader evaluation, rather than its novelty; otherwise [...] the revised algorithm should be included explaining what the differences with the Algorithm 2 in the conference paper are.}
\resp We thank the reviewer for requiring this clarification. [To be confirmed against the conference sources. If the classification algorithm is unchanged, the contribution is rephrased as expanded evaluation and deeper type-aware connectivity treatment, and if it was refined, the delta will be stated explicitly with the revised algorithm included.] The revised contribution list already no longer presents classification itself as novel.
\loc{Section 1 (Introduction) and Section 7.2 (Door Classification) [p.~XX]. Marked in purple.}

\commentthree{R3.1c}{[...] there are several discrepancies that need to be clarified, e.g., total number of rooms (i.e., 22 vs 3x9=27), resolution (i.e., Variable vs 2500x1600), etc. Importantly, authors should include also detailed results from at least one more multi-floor building, i.e., `site B' [...] I would expect to see a baseline comparison with at least one of the competing indoor mapping approaches listed in Table 1 [...]}
\resp We thank the reviewer for all three demands, which we address in turn. First, the dataset characteristics are reconciled in the revised dataset table, which lists per-plan sizes and counts explicitly. [The remaining cross-checks against the conference table will be completed with the consolidated dataset.] Second, Site B is being restored to the evaluation as described under R1.6. Third, on baselines, the revision adds a quantitative comparison against a structure-only medial-axis navigation graph, a classical alternative that we deliberately strengthen by granting it our own room and exit detections so the comparison isolates routing structure. The outcome is decisive. The skeleton needs between 3.7 and 12 times more nodes, fails completeness at 0.78 to 0.83 because it fragments at doorway gaps, produces worse route stretch against the same geodesic reference, answers queries between 37 and 94 times slower, and expresses none of the typed downstream tasks. [Whether a published learned floorplan-parsing system can additionally be run under comparable conditions is under investigation, and if it proves infeasible the revision will provide a structured qualitative comparison in its place.]
\loc{Table 2 in Section 10.1 [p.~XX] and Section 10.7 (Baseline Comparison) with Table 4 [p.~XX]. Marked in purple.}

\commentthree{R3.2}{The degree of automation is somewhat overstated for the multi-floor setting. [...] inter-floor transition mappings are provided as structured input and then validated by the system. This is still useful, but it is not the same as fully automatic inference of vertical correspondences from the floor plans alone. I recommend making this limitation more explicit and tempering claims of end-to-end automation accordingly.}
\resp We thank the reviewer for the measured critique and chose the stronger of the two available paths. Instead of tempering the claims to match the system, we changed the system to match the claims. Vertical correspondences are now inferred automatically from the floor plans through registration and assignment with rejection, described under R1.9, with hand-written mappings retained as an override. The remaining metadata assistance, floor indices from filenames, is stated explicitly rather than folded into the word automatic, and every automation claim in the abstract, introduction, and Section 9 has been re-audited against the implementation.
\loc{Section 9.2 [p.~XX], plus the claim audit in the Abstract and Section 1 [p.~XX]. Marked in vermilion, jointly with R1.9 and R2.W1.}

\commentthree{R3.3}{The main failure mode, i.e., corridor-to-corridor door detection, has a disproportionate effect on connectivity, but the mitigation strategy is weak. [...] I suggest adding at least one of the following: (a) a stronger model or heuristic fallback for corridor connectors, (b) an ablation showing how much completeness recovers under modest detector improvements, beyond the current high-level projection, or (c) some topological repair strategy [...]}
\resp We thank the reviewer for offering three concrete routes forward. The revision delivers option (b) in full and lays the measured groundwork for option (c). The high-level projection is replaced by Monte-Carlo knockout and recovery curves with error bars, random and by type, plus per-door sensitivity analysis. The study shows the damage concentrates in bridge corridor doors, where a single door costs 0.51 of completeness on its own, while most corridor doors are mesh-redundant. Because bridge doors coincide with the graph's cut vertices, they are mechanically identifiable, which converts the failure mode into a targeted verification and repair problem. The discussion states this mechanism and scopes full image-guided repair as immediate future work rather than claiming it prematurely. Option (a), a stronger detector, is explicitly declared out of scope for this revision.
\loc{Section 10.3 (Connectivity under Door-Detector Degradation) with Figure 3 [p.~XX] and Section 11 (Discussion, Door Detection and Repair) [p.~XX]. Marked in vermilion, jointly with R2.W3.}

\commentthree{R3.4}{The claims around general applicability should be narrowed to match the actual input assumptions. [...] These assumptions are reasonable, but they exclude many low-quality, style-diverse, or sparsely annotated floor plans. The paper should state more clearly where the system is expected to work reliably and where it is not.}
\resp We thank the reviewer for the framing, which we adopt directly as an operating envelope. The revised input section gives the required label vocabulary as a table with per-class requirements, states the annotation requirement as a hard precondition, and notes wall recoverability by thresholding and the door-symbol convention dependency. It then states where the system is not expected to work, namely hand-drawn sketches, sparsely annotated plans, and non-standard symbol sets, and it documents the failure behavior on a missing label class, rejection rather than silent mis-segmentation, with a demonstrated example.
\loc{Section 3.1 (Input Assumptions) with Table 1 [p.~XX] and Section 11 (Discussion, Operating Envelope) [p.~XX]. Marked in purple.}

\commentthree{R3.5}{The geometric-fidelity analysis (section 10.7) is useful but incomplete as a validation of navigational realism. [...] it would strengthen the paper to compare against manually verified route distances or some ground-truth navigational paths for a subset of room pairs, especially for cross-floor movement where nominal transition weights can distort fidelity.}
\resp We thank the reviewer for both halves of this comment, and each produced a change. For fidelity, the revision adopts the free-space geodesic reference with visual verification overlays, described under R1.7. For the sharp observation about nominal transition weights, which was entirely correct, inter-floor edges no longer carry a nominal cost. They are weighted by modeled vertical travel, the product of a travel factor, the story height, and the drawing scale in pixels per meter, where the scale is bootstrapped from detected door widths under the standardized 0.9-meter leaf. The estimates are independently consistent across floors at 31.1 versus 32.2 pixels per meter, which supports the bootstrap, and the correction measurably matters, since the step-free detour ratio corrects from 2.02 to 1.79. The same scale estimate also feeds inter-floor registration as a prior.
\loc{Section 9.3 (Inter-Floor Edge Weights) [p.~XX], Section 10.6 for fidelity [p.~XX], and the corrected detour figures in Sections 10.8 and 10.9 [p.~XX]. Marked in vermilion for the shared fidelity work and purple for the weight model.}

\bigskip
\noindent We hope the revision demonstrates how seriously we took each point. The reviews improved the system itself, not only its presentation, and we are grateful for that. We remain happy to provide any further clarification the reviewers or editors may wish.

\bigskip
\noindent Sincerely,\\
The Authors

\end{document}
